\documentclass[10pt,a4paper]{amsart}
\usepackage{amsmath,amssymb,amsthm,graphicx}
\usepackage{verbatim}
\newenvironment{reference}{\comment}{\endcomment}
%\newenvironment{reference}{}{}
\newenvironment{slogan}{\comment}{\endcomment}
\newenvironment{history}{\comment}{\endcomment}
% For commutative diagrams we use Xy-pic
\usepackage[all]{xy}

% We use 2cell for 2-commutative diagrams.
\xyoption{2cell}
% Theorem environments.
%
\theoremstyle{plain}
\newtheorem{theorem}[subsection]{Theorem}
\newtheorem{proposition}[subsection]{Proposition}
\newtheorem{lemma}[subsection]{Lemma}

\theoremstyle{definition}
\newtheorem{definition}[subsection]{Definition}
\newtheorem{example}[subsection]{Example}
\newtheorem{exercise}[subsection]{Exercise}
\newtheorem{situation}[subsection]{Situation}

\theoremstyle{remark}
\newtheorem{remark}[subsection]{Remark}
\newtheorem{remarks}[subsection]{Remarks}

\numberwithin{equation}{subsection}

% Macros
%
\def\lim{\mathop{\mathrm{lim}}\nolimits}
\def\colim{\mathop{\mathrm{colim}}\nolimits}
\def\Spec{\mathop{\mathrm{Spec}}}
\def\Hom{\mathop{\mathrm{Hom}}\nolimits}
\def\Ext{\mathop{\mathrm{Ext}}\nolimits}
\def\SheafHom{\mathop{\mathcal{H}\!\mathit{om}}\nolimits}
\def\SheafExt{\mathop{\mathcal{E}\!\mathit{xt}}\nolimits}
\def\Sch{\mathit{Sch}}
\def\Mor{\mathop{\mathrm{Mor}}\nolimits}
\def\Ob{\mathop{\mathrm{Ob}}\nolimits}
\def\Sh{\mathop{\mathit{Sh}}\nolimits}
\def\NL{\mathop{N\!L}\nolimits}
\def\CH{\mathop{\mathrm{CH}}\nolimits}
\def\proetale{{pro\text{-}\acute{e}tale}}
\def\etale{{\acute{e}tale}}
\def\QCoh{\mathit{QCoh}}
\def\Ker{\mathop{\mathrm{Ker}}}
\def\Im{\mathop{\mathrm{Im}}}
\def\Coker{\mathop{\mathrm{Coker}}}
\def\Coim{\mathop{\mathrm{Coim}}}

% Boxtimes
%
\DeclareMathSymbol{\boxtimes}{\mathbin}{AMSa}{"02}

%
% Macros for moduli stacks/spaces
%
\def\QCohstack{\mathcal{QC}\!\mathit{oh}}
\def\Cohstack{\mathcal{C}\!\mathit{oh}}
\def\Spacesstack{\mathcal{S}\!\mathit{paces}}
\def\Quotfunctor{\mathrm{Quot}}
\def\Hilbfunctor{\mathrm{Hilb}}
\def\Curvesstack{\mathcal{C}\!\mathit{urves}}
\def\Polarizedstack{\mathcal{P}\!\mathit{olarized}}
\def\Complexesstack{\mathcal{C}\!\mathit{omplexes}}
% \Pic is the operator that assigns to X its picard group, usage \Pic(X)
% \Picardstack_{X/B} denotes the Picard stack of X over B
% \Picardfunctor_{X/B} denotes the Picard functor of X over B
\def\Pic{\mathop{\mathrm{Pic}}\nolimits}
\def\Picardstack{\mathcal{P}\!\mathit{ic}}
\def\Picardfunctor{\mathrm{Pic}}
\def\Deformationcategory{\mathcal{D}\!\mathit{ef}}
\usepackage{cleveref}
\crefname{originallemma}{original lemma}{original lemmas}
\Crefname{originallemma}{Original lemma}{Original lemmas}
\makeatletter
\newlabel{lemma-diagonal-diagonal}{{lemma-diagonal-diagonal}{source}{}{sourcepointer-1}{}}
\newlabel{lemma-diagonal-diagonal@cref}{{[originallemma][0][]lemma-diagonal-diagonal}{[1][0][]source}{}{}{}}
\makeatother
\crefname{originalreference}{original reference}{original references}
\Crefname{originalreference}{Original reference}{Original references}
\makeatletter
\newlabel{spaces-groupoids-lemma-open-over-which-unramified-or-lqf}{{spaces-groupoids-lemma-open-over-which-unramified-or-lqf}{source}{}{sourcepointer-2}{}}
\newlabel{spaces-groupoids-lemma-open-over-which-unramified-or-lqf@cref}{{[originalreference][0][]spaces-groupoids-lemma-open-over-which-unramified-or-lqf}{[1][0][]source}{}{}{}}
\makeatother
\crefname{originalreference}{original reference}{original references}
\Crefname{originalreference}{Original reference}{Original references}
\makeatletter
\newlabel{stacks-properties-lemma-substacks-presentation}{{stacks-properties-lemma-substacks-presentation}{source}{}{sourcepointer-3}{}}
\newlabel{stacks-properties-lemma-substacks-presentation@cref}{{[originalreference][0][]stacks-properties-lemma-substacks-presentation}{[1][0][]source}{}{}{}}
\makeatother
\title{A point lies in the quasi-DM locus exactly when its automorphism group is locally quasi-finite}
\author{The Stacks Project Authors}
\date{stacks-project:stacks-morphisms.tex:local-snapshot-20261012; GNU FDL 1.2 or later}
\begin{document}
\maketitle
\noindent\textit{Editorial note.} All delivered mathematical text is the author's own, copied from the pinned Stacks Project chapter source sources/stacks-algebraic-stacks-morphisms/pinned-original-source.tex; no statement or proof step was rewritten, reordered or removed.\par\medskip
\noindent\textit{Editorial note.} Preamble provenance: the Stacks Project keeps one single shared preamble for every chapter; this card takes it directly from the staged sources/stacks-algebraic-stacks-morphisms/preamble.tex as three verbatim spans (lines 8-12, 14-18 and 154-217), so no derived file was written and nothing was added or redefined. The wrapper class and the base packages are supplied by the assembler.\par\medskip
\noindent\textit{Editorial note.} Bibliography provenance: neither the counted statement nor its proof contains a citation (this lane rescans the whole span before writing the card), so the card carries no bibliography and no entry was reconstructed.\par\medskip
\noindent\textit{Editorial note.} Selection provenance: this card was opened under the file-based control contract (CONTROL.md), with the mechanical pre-card gate run on its statement and proof interval before the card was written.\par\medskip
\noindent\textit{Allowed standard prerequisites (editorial).}\begin{itemize}
\item The existence of a maximal open subscheme over which the automorphism group space is locally quasi-finite, with base change formation. Invoked at the start of the counted proof (source lines 4832-4836) as a cross-chapter pointer. Produces the open subscheme U' of the presentation whose invariant descent gives the quasi-DM locus.
\item The correspondence between open substacks of a presentation and invariant open subschemes. Invoked inside the counted proof (source lines 4841-4844) as a cross-chapter pointer. Descends the open subscheme U' to an open substack of the stack.
\item The characterization of the quasi-DM locus through the diagonal being locally quasi-finite. Invoked inside the counted proof (source lines 4849-4851). Turns local quasi-finiteness of the inertia over the descended substack into quasi-DM-ness.
\end{itemize}\medskip
\ifdefined\hypertarget\hypertarget{sourcepointer-1}{}\fi
\noindent\textit{Source reference.} Original lemma lemma-diagonal-diagonal: sources/stacks-algebraic-stacks-morphisms/pinned-original-source.tex, source lines 3959--3959. Navigation pointer only: the counted text refers to the Stacks item lemma-diagonal-diagonal, which belongs to another chapter of the project and is not part of this card. The pointer carries the target's own label as its number because this chapter's staged snapshot does not contain the other chapter's tag table; its content is not reproduced here.\par\medskip
\ifdefined\hypertarget\hypertarget{sourcepointer-2}{}\fi
\noindent\textit{Source reference.} Original reference spaces-groupoids-lemma-open-over-which-unramified-or-lqf: sources/stacks-algebraic-stacks-morphisms/pinned-original-source.tex, source lines 4840--4840. Navigation pointer only: the counted text refers to the Stacks item spaces-groupoids-lemma-open-over-which-unramified-or-lqf, which belongs to another chapter of the project and is not part of this card. The pointer carries the target's own label as its number because this chapter's staged snapshot does not contain the other chapter's tag table; its content is not reproduced here.\par\medskip
\ifdefined\hypertarget\hypertarget{sourcepointer-3}{}\fi
\noindent\textit{Source reference.} Original reference stacks-properties-lemma-substacks-presentation: sources/stacks-algebraic-stacks-morphisms/pinned-original-source.tex, source lines 4851--4851. Navigation pointer only: the counted text refers to the Stacks item stacks-properties-lemma-substacks-presentation, which belongs to another chapter of the project and is not part of this card. The pointer carries the target's own label as its number because this chapter's staged snapshot does not contain the other chapter's tag table; its content is not reproduced here.\par\medskip
\begin{lemma}
\label{lemma-points-DM-locus}
Let $\mathcal{X}$ be an algebraic stack. Let $x \in |\mathcal{X}|$
correspond to $x : \Spec(k) \to \mathcal{X}$. Let $G_x/k$
be the automorphism group algebraic space of $x$. Then
\begin{enumerate}
\item  $x$ is in the DM locus of $\mathcal{X}$
if and only if $G_x \to \Spec(k)$ is unramified, and
\item $x$ is in the quasi-DM locus of $\mathcal{X}$
if and only if $G_x \to \Spec(k)$ is locally quasi-finite.
\end{enumerate}
\end{lemma}

\begin{proof}
Proof of (2). Choose a scheme $U$ and a surjective smooth morphism
$U \to \mathcal{X}$. Consider the fibre product
$$
\xymatrix{
G \ar[r] \ar[d] & \mathcal{I}_\mathcal{X} \ar[d] \\
U \ar[r] & \mathcal{X}
}
$$
Recall that $G$ is the automorphism group algebraic space of
$U \to \mathcal{X}$. By Groupoids in Spaces, Lemma
\ref{spaces-groupoids-lemma-open-over-which-unramified-or-lqf}
there is a maximal open subscheme $U' \subset U$
such that $G_{U'} \to U'$ is locally quasi-finite.
Moreover, formation of $U'$ commutes with arbitrary
base change. In particular the two inverse images of $U'$
in $R = U \times_\mathcal{X} U$ are the same open subspace of $R$
(since after all the two maps $R \to \mathcal{X}$ are isomorphic
and hence have isomorphic automorphism group spaces).
Hence $U'$ is the inverse image of an open substack
$\mathcal{X}' \subset \mathcal{X}$ by
Properties of Stacks, Lemma
\ref{stacks-properties-lemma-substacks-presentation}
and we have a cartesian diagram
$$
\xymatrix{
G_{U'} \ar[r] \ar[d] & \mathcal{I}_{\mathcal{X}'} \ar[d] \\
U' \ar[r] & \mathcal{X}'
}
$$
Thus the morphism $\mathcal{I}_{\mathcal{X}'} \to \mathcal{X}'$
is locally quasi-finite and we conclude that
$\mathcal{X}'$ is quasi-DM by Lemma \ref{lemma-diagonal-diagonal}
part (5). On the other hand, if $\mathcal{W} \subset \mathcal{X}$
is an open substack which is quasi-DM, then the inverse image
$W \subset U$ of $\mathcal{W}$ must be contained in $U'$ by our
construction of $U'$ since
$\mathcal{I}_\mathcal{W} =
\mathcal{W} \times_\mathcal{X} \mathcal{I}_\mathcal{X}$
is locally quasi-finite over $\mathcal{W}$.
Thus $\mathcal{X}'$ is the quasi-DM locus.
Finally, choose a field extension $K/k$ and a $2$-commutative
diagram
$$
\xymatrix{
\Spec(K) \ar[r] \ar[d] & \Spec(k) \ar[d]^x \\
U \ar[r] & \mathcal{X}
}
$$
Then we find an isomorphism
$G_x \times_{\Spec(k)} \Spec(K) \cong G \times_U \Spec(K)$
of group algebraic spaces over $K$. Hence $G_x$ is locally quasi-finite
over $k$ if and only if $\Spec(K) \to U$ maps into $U'$
(use the commutation of formation of $U'$ and
Groupoids in Spaces, Lemma
\ref{spaces-groupoids-lemma-open-over-which-unramified-or-lqf}
applied to $\Spec(K) \to \Spec(k)$ and $G_x$ to see this).
This finishes the proof of (2). The proof of (1) is
exactly the same.
\end{proof}

\end{document}
